\begin{definition}[The actual correlated two-bond charge insertion]
    \label{def:peps_regular_correlated_charge_column}
    \lean{TNLean.PEPS.regularProjectorChargePairOpenRegionMatrix,
      TNLean.PEPS.regularChargePairOpenRegionMatrix}
    \leanok
    \uses{def:peps_regular_charge_pair, def:peps_regular_projector_open_region}
    Let $R$ be a finite region of an ordered simple graph, let $e_0,e_1$
    be internal bonds, and let $A_v$ be the original site coefficients.
    For a physical configuration $s$ and crossing labels $\theta$, put
    \[
        C_{\chi,p}(s,\theta)
          =\sum_{\eta|_{\partial R}=\theta}
             \chi(p\eta_{e_1}^{-1}\eta_{e_0})
             \prod_{v\in R} A_v(\eta|_v,s_v).
    \]
    Replacing $A_v$ by the local regular averaging projectors gives the
    canonical column $B_{\chi,p}$. This is a single correlated weight on
    two bonds, as in \cite[charge-pair creation]{Schuch2010PEPS},
    lines 2505--2558.
\end{definition}

\begin{theorem}[The correlated weight in actual tree coordinates]
    \label{thm:peps_regular_correlated_charge_coordinates}
    \lean{TNLean.PEPS.regularProjectorChargePairOpenRegionMatrix_coordinates,
      TNLean.PEPS.regularChargePairOpenRegionMatrix_eq_image,
      TNLean.PEPS.regularChargePairOpenRegionMatrix_one_eq_openRegionWeight}
    \leanok
    \uses{def:peps_regular_correlated_charge_column,
      thm:peps_regular_region_projector_coordinates}
    Choose a spanning tree of the induced region and a root. In its actual
    coordinates $c=(y,a,r,z)$,
    \[
        B_{\chi,p}(c,\theta)
           =\chi(pa_{e_1}^{-1}a_{e_0}) B_{1,1}(c,\theta).
    \]
    If every original site is regular $G$-injective, its product physical
    map $T_R$ satisfies $T_R B_{\chi,p}=C_{\chi,p}$. In particular,
    $C_{1,1}$ is the original open-region contraction.
\end{theorem}
\begin{proof}\leanok
    A nonzero product of local projectors supplies compatible vertex
    translations. Tree synchronization writes them as $q_v=r_v x$.
    Thus $a_e=x\eta_e$, and the common $x$ cancels from the relative
    coordinate. The physical identity follows from $T_vP_v=T_v$ at each
    site and the literal contraction sum.
\end{proof}

\begin{definition}[A unitary on two internal reference labels]
    \label{def:peps_regular_charge_reference_matrix}
    \lean{TNLean.PEPS.regularRegionChargeReferenceMatrix}
    \leanok
    \uses{thm:peps_regular_correlated_charge_coordinates}
    For distinct internal bonds $e_0,e_1$, split the two references
    $(a_{e_0},a_{e_1})$ from all other tree coordinates. A matrix $Q$ on
    these two group registers acts as $Q\otimes I$ in this decomposition.
    Transport this operation through the actual physical coordinate
    permutation to obtain $U_R(Q)$.
\end{definition}

\begin{theorem}[Boundary-uniform preparation on the original spins]
    \label{thm:peps_regular_physical_charge_pair_creation}
    \lean{
      TNLean.PEPS.exists_unitary_regularChargeReferencePreparation,
      TNLean.PEPS.regularRegionChargeReferenceMatrix_mem_unitaryGroup,
      TNLean.PEPS.regularRegionChargeReferenceMatrix_commute_vertexTranslation,
      TNLean.PEPS.regularRegionChargeReferenceMatrix_commute_localProjector,
      TNLean.PEPS.regularRegionChargeReferenceMatrix_mulVec,
      TNLean.PEPS.exists_unitary_regularPhysicalChargePairCreation}
    \leanok
    \uses{def:peps_regular_charge_reference_matrix,
      thm:peps_regular_charge_pair_normalization,
      thm:peps_regular_charge_labels,
      lem:peps_unitary_vector_swap,lem:peps_permutation_matrix_commutation}
    For a unitary irreducible character $\chi$ and $p\in G$, there is a
    unitary $Q$ commuting with simultaneous left translation such that
    \[
        Q\bigl(|G|^{-1}\mathbf1\bigr)(a,b)
           =|G|^{-1}\chi(pb^{-1}a).
    \]
    The resulting $U_R(Q)$ is unitary and commutes with every independent
    vertex translation, hence with the product of local averaging projectors.
    It sends $B_{1,1}(\cdot,\theta)$ to $B_{\chi,p}(\cdot,\theta)$ for every
    $\theta$.

    Suppose now that the original sites are regular $G$-isometric and
    that $\chi$ is an actual irreducible character of the regular
    representation. There is one original-spin unitary $W$, chosen before
    all crossing configurations, such that
    \[
        W C_{1,1}(\cdot,\theta)=C_{\chi,p}(\cdot,\theta)
        \qquad\text{for every }\theta.
    \]
    This is an auxiliary finite-region realization of the construction in
    \cite[charge-pair creation]{Schuch2010PEPS}, lines 2505--2558.
\end{theorem}
\begin{proof}\leanok
    The normalized uniform and charge vectors have equal norm. The identity
    handles the trivial character; otherwise their orthogonality supplies
    a unitary exchange. Both vectors are fixed by simultaneous translation,
    so this exchange commutes with that action. Invariance of matrix entries
    under a permutation implies commutation with its permutation matrix.
    The actual tree-coordinate transformation then proves commutation with
    every vertex action. Finally, the local isometry assumptions derive
    $T_R^*T_R=c_RP_R$ with $c_R>0$. The resulting supported unitary extends
    to a unitary on the original physical region, with $WT_R=T_RU_R(Q)$.
    Apply this equality to the actual canonical boundary columns.
\end{proof}

\begin{theorem}[Actual six-spin charge-pair preparation]
    \label{thm:peps_torus_physical_charge_pair_creation}
    \lean{TNLean.PEPS.IsGIsometric.exists_unitary_torusPhysicalChargePairCreation}
    \leanok
    \uses{thm:peps_regular_physical_charge_pair_creation}
    Let a finite torus have horizontal period at least three and vertical
    period at least four. Let its original four-leg tensor be regular
    $G$-isometric. At every translated two-column, three-row block,
    the four vertical bonds and the outer horizontal bond form the
    five-edge tree. The two remaining horizontal chords are $e_0,e_1$.
    For every regular irreducible character $\chi$ and $p\in G$, one unitary
    on these six original spins sends the unmodified open contraction to
    $C_{\chi,p}$, uniformly in all crossing labels.

    The two horizontal chords are the top and middle references in the
    source diagram after reversing its vertical drawing axis. Native bond
    labels and sorting are retained even when the block crosses a seam.
    This proves the stated finite-torus regular specialization of the
    six-spin preparation; it makes no parent-Hamiltonian or ground-space
    identification.
\end{theorem}
\begin{proof}\leanok
    The translated induced block supplies its tree, root and two distinct
    chords. Transport the original four-leg local isometry to its incident
    coordinates and apply the preceding preparation theorem. The literal
    unweighted column is the actual original open-region contraction.
\end{proof}
